% !TEX root = ../main.tex
\begin{abstract}
Timber harvest schedules computed by linear programming are solved once, as if every future decision could be fixed today, but planners re-solve every period. In his 1991 doctoral thesis, P. J. Daugherty showed that such open-loop plans can be dynamically inconsistent: later planners, re-solving a model of the same form, would not follow them. We provide an open, reproducible benchmark of his analysis, built from public Umpqua National Forest planning documents. Across 432 scenarios (18 initial forests, four discount rates, six harvest-flow policies), scored over the thesis's observation window, inconsistency occurred in 49\% of scenarios with harvest-flow constraints, at smaller magnitudes than the thesis reported; without such constraints, it occurred only at discount rates of 0--2\% and only when each replan extended the horizon. On the main all-mature forest, the realized harvest stayed below the promised non-declining level throughout. Inconsistency nearly vanished when each replan kept the original end date and the flow constraint on the previous harvest. In mitigation experiments, declining discount rates re-applied by each planner mostly increased inconsistency, a revenue-denominated flow constraint did not reduce it, a realized-harvest-anchored constraint did, and a calibrated harvest cap proposed in the thesis nearly removed it, largely by construction.
\end{abstract}
